\section{Fundamentals}
The properties of the ultimate display follow from human vision. Before discussing display technology, it is therefore necessary to understand exactly why a two-dimensional image of an object looks less convincing than the real thing.
\subsection{Depth perception}
Depth perception is typically analyzed in terms of various contributing factors called depth cues.
\begin{center}
\begin{tabular}{l p{2.5cm} p{2.5cm}}
 & \textbf{Monocular} & \textbf{Binocular} \\ 
 \hline
 \textbf{Visual} & Pictorial cues & Retinal disparity \\  
 \hline
 \textbf{Oculomotor} & Accomodation, Pupillary \par response & Convergence \\
\end{tabular}
\end{center}
Retinal disparity arises from the slight difference in perspective between the left and the right eye. Light originating from the same point in space does not necessarily reach corresponding points on the left and right retinas. A detailed explanation of the geometry is given by Patterson and Martin (1992). \cite{patterson1992human} The oculomotor cues are sensed by proprioceptors in the muscles responsible for focusing the eye lens (Accomodation), turning the eyeballs (Convergence), and constricting as well as dilating the pupil in response to brightness (Pupillary response). The way the eyes adjust in these three ways when viewing an object correlates to its distance. This allows the visual system to use this form of proprioception to perceive depth.
\subsection{The plenoptic function}
\subsection{Light fields}